% ==================== 第六章 ====================
\chapter{制度护航：治理体系与国际博弈}

\textbf{本章聚焦"制度"建设——治理体系的设计和国际话语权的构建。}

技术能力必须有制度保障才能转化为持续竞争力。本章提供治理体系的具体制度设计、国际博弈的具体策略和话语框架、以及信息安全管理的创新机制。

本章说明：人才问题因其战略重要性，将在第九章"人才强基"中单独系统论述，包括"人才安全"理论框架、人才资产负债表模型、以及"磁吸-锚定-生态"三位一体战略。本章仅概述人才问题的紧迫性，作为治理体系建设的背景。

\section{人才危机：治理体系必须回应的紧迫挑战}

在展开治理体系设计之前，我们必须首先认识到一个严峻的现实：我国AI人才，特别是AI安全人才，面临着超过90\%的缺口。国际对标数据表明，Anthropic对齐团队约150人、OpenAI安全团队约100人、DeepMind安全团队约200人，而我国真正在做前沿对齐研究的不超过30人，能做系统性模型安全评估的红队测试人员不超过100人。

更令人担忧的是人才流失问题。2024-2025年间，多位国内顶尖AI研究者加入海外实验室。薪酬差距（顶尖研究者美国年薪可达100-300万美元）、研究环境差距、评价体系问题和行政负担，都是导致流失的重要原因。这一人才危机是治理体系设计必须正视和回应的核心挑战。关于人才战略的系统性分析和政策建议，详见第九章。

\section{治理体系：从被动应对到主动设计}

\subsection{核心问题：谁来管、怎么管、管什么}

当前的治理困境在于部门分割、缺乏统一协调。科技部管基础研究，工信部管产业发展，网信办管内容安全，公安部管安全应用。各部门各管一段，在AI这样一个需要系统性治理的领域，这种分割式管理模式难以形成合力。

国际比较为我们提供了有益参考。美国由白宫科技政策办公室（OSTP）与国家安全委员会（NSC）协调，各部门分工执行。英国由科学、创新与技术部（DSIT）统一负责AI政策。欧盟则由欧委会AI办公室统一协调。这些模式的共同点是在顶层建立统一的协调机制。

\subsection{治理架构设计}

建议的治理架构包括"一委、一办、多中心"三个层次。

"一委"即国家AI发展与安全委员会，应设在中央层面，由政治局常委分管，负责顶层设计、重大决策和跨部门协调。成员包括相关部委主要负责人，每季度召开一次全体会议，重大事项随时开会研究。

"一办"即国家AI治理办公室，可定为正部级或副部级，设在科技部或单独设立。主要职能包括日常协调、政策制定、标准管理和国际合作，编制规模约100-200人。

"多中心"指专业支撑机构，包括：国家AI安全测评中心负责模型安全评估和认证，国家AI标准研究中心负责标准制定和国际对接，国家AI伦理研究中心负责伦理审查和指南制定，国家AI应急响应中心负责安全事件应急处置。

\subsection{法规体系完善}

\textbf{立法优先序}：

\begin{table}[H]
\centering
\caption{AI立法路线图}
\small
\begin{tabular}{p{1.5cm}p{3cm}p{4cm}p{3.5cm}}
\toprule
\textbf{优先级} & \textbf{法规名称} & \textbf{核心内容} & \textbf{建议时间} \\
\midrule
P0 & 《人工智能法》 & AI发展与治理基本法 & 2025年立项，2027年出台 \\
P1 & 《AI安全管理条例》 & AI系统安全要求 & 2025年出台 \\
P1 & 《AI算力管理办法》 & 智算中心管理规范 & 2025年出台 \\
P2 & 《AI数据管理条例》 & 训练数据合规要求 & 2026年出台 \\
P2 & 《AI应用分级管理办法》 & 高风险应用准入 & 2026年出台 \\
\bottomrule
\end{tabular}
\end{table}

\subsection{信息安全管理优化}

传统保密制度不适应AI时代的信息聚合能力，这是一个必须正视的新问题。

"反马赛克"审查机制的触发条件包括：涉及领导人行程、人事任免、重大项目等敏感信息；跨部门信息同时发布；以及可能被AI聚合推断的公开信息。

审查流程分为四个步骤。第一步是自动预筛，AI系统自动检测信息的敏感度评分。第二步是模拟推演，用多个主流大模型尝试聚合推断。第三步是人工复核，安全专家评估推断结果的风险。第四步是处置决定，根据评估结果选择延迟发布、脱敏发布或正常发布。

组织保障方面，需要在国家保密局设立"AI安全审查处"，在各部委指定信息发布审查专员，并配备AI审查工具和专业培训。时限要求为：常规审查24小时内完成，紧急审查4小时内完成，申诉复核72小时内完成。

\section{国际博弈：竞争中寻求合作空间}

%% =============== 新增：数据主权与隐私计算技术详解 ===============

\section{数据主权与数字边疆：技术实现方案}

\subsection{数据主权的技术内涵}

数据主权不仅是法律概念，更需要技术手段保障。本节详细分析数据主权的技术实现路径\cite{DataSovereignty2024, EUDataAct2024, ChinaDataSecurity2024}。

数据主权的技术实现分三个层次。第一层是数据本地化：关键数据必须存储在境内服务器，数据处理必须在境内完成，境外镜像需符合安全评估要求。第二层是跨境流动管控：重要数据出境需经过安全评估，跨境传输必须端到端加密，禁止向特定国家或地区传输敏感数据。第三层是主权可控计算：通过联邦学习实现"数据不动模型动"，确保数据在指定条件下被彻底销毁（可验证删除），全程追踪数据使用情况（使用审计）。

\subsection{深度技术详解：联邦学习架构}

联邦学习（Federated Learning）是实现"数据不出域"的核心技术\cite{FederatedLearning2024, GoogleFL2024, FedAvg2024}。

\textbf{1. 基础FedAvg算法}

每轮通信$t$，参与方$k$本地更新：
\begin{equation}
w_k^{t+1} = w_k^t - \eta \nabla F_k(w_k^t)
\end{equation}

服务器聚合：
\begin{equation}
w^{t+1} = \sum_{k=1}^{K} \frac{n_k}{n} w_k^{t+1}
\end{equation}

其中$n_k$为第$k$方数据量，$n=\sum_k n_k$为总数据量。

\textbf{2. 异构数据下的改进：FedProx}

添加近端正则项解决Non-IID数据问题：
\begin{equation}
\min_{w_k} F_k(w_k) + \frac{\mu}{2} \|w_k - w^t\|^2
\end{equation}

正则系数$\mu$控制本地模型与全局模型的距离。

\textbf{3. 个性化联邦学习：APFL}

每方维护全局模型$w$和本地模型$v_k$：
\begin{equation}
\bar{w}_k = \alpha_k v_k + (1-\alpha_k) w
\end{equation}

混合系数$\alpha_k$根据本地数据分布自适应调整。

\textbf{4. 通信效率优化}
梯度压缩：Top-K稀疏化，仅传输最大的K个梯度分量；量化通信：将FP32梯度量化为INT8甚至1-bit；周期性聚合：减少通信频率，增加本地迭代次数。

\textbf{攻击类型}：

\textbf{1. 梯度逆向攻击}
原理：从梯度重建训练数据；方法：通过优化 $\min_x \|g - \nabla_w L(w, x)\|^2$ 来恢复原始样本$x$；效果：可恢复图像、文本等敏感数据。

\textbf{2. 模型投毒攻击}
原理：恶意参与方上传被污染的模型更新；目标：使全局模型在特定输入上输出错误结果；后门攻击：植入触发器，正常输入正常输出，特定触发输入恶意输出。

\textbf{3. 成员推断攻击}
原理：判断某条数据是否参与了训练 方法：利用模型对训练数据和非训练数据的不同行为

\textbf{防御方案}：

\textbf{差分隐私（Differential Privacy）}
\begin{equation}
\Pr[\mathcal{M}(D) \in S] \leq e^\varepsilon \cdot \Pr[\mathcal{M}(D') \in S] + \delta
\end{equation}

其中$D, D'$为相差一条记录的数据集，$(\varepsilon, \delta)$-DP保证个体隐私。

具体实现：梯度裁剪+高斯噪声
\begin{equation}
\tilde{g} = \frac{g}{\max(1, \|g\|/C)} + \mathcal{N}(0, \sigma^2 C^2 \mathbf{I})
\end{equation}

$C$为裁剪阈值，$\sigma$为噪声系数，典型配置$\sigma=1.0$, $C=1.0$可达$\varepsilon=3$。

\textbf{安全聚合（Secure Aggregation）}
原理：服务器只能看到聚合结果，看不到个体更新；方法：基于秘密分享或同态加密；开销：通信量增加$O(K^2)$，计算量增加$O(K)$。

\subsection{深度技术详解：同态加密}

同态加密（Homomorphic Encryption）允许在密文上直接计算\cite{HomomorphicEncryption2024, TFHE2024, CKKS2024}。

\textbf{1. 同态加密分类}

\begin{table}[H]
\centering
\small
\begin{tabular}{lccc}
\toprule
\textbf{类型} & \textbf{支持运算} & \textbf{代表方案} & \textbf{应用场景} \\
\midrule
部分同态(PHE) & 仅加法或乘法 & Paillier, RSA & 简单统计 \\
有限同态(SWHE) & 有限次乘法 & BGV, BFV & 浅层神经网络 \\
全同态(FHE) & 任意次数 & TFHE, CKKS & 深度学习 \\
\bottomrule
\end{tabular}
\end{table}

\textbf{2. CKKS方案详解（近似计算最优）}

CKKS适合机器学习中的浮点运算：
明文空间：$\mathbb{C}^{N/2}$复数向量；编码：将浮点向量编码为多项式；加法同态：$\text{Enc}(m_1) + \text{Enc}(m_2) = \text{Enc}(m_1 + m_2)$；乘法同态：$\text{Enc}(m_1) \times \text{Enc}(m_2) = \text{Enc}(m_1 \cdot m_2)$。

关键参数：
多项式度数$N$：典型$N=2^{15}$或$2^{16}$，决定批处理槽位数；模数链$\{q_0, q_1, ..., q_L\}$：$L$决定乘法深度；缩放因子$\Delta$：控制精度，典型$\Delta=2^{40}$。

\textbf{3. Bootstrapping技术}

FHE的核心挑战是噪声增长。Bootstrapping可刷新密文：
\begin{equation}
\text{Bootstrap}(\text{Enc}(m)) \rightarrow \text{Enc}'(m)
\end{equation}

代价：
TFHE Bootstrapping：约10-50ms/门；CKKS Bootstrapping：约1-10秒/次；是FHE实用化的主要瓶颈。

\textbf{4. FHE神经网络推理性能}\footnote{数据来源：综合学术文献，包括Microsoft SEAL、IBM HELib等项目的基准测试。实际性能因硬件和参数设置差异较大。}

\begin{table}[H]
\centering
\small
\begin{tabular}{lccc}
\toprule
\textbf{模型} & \textbf{明文延迟} & \textbf{FHE延迟} & \textbf{放大倍数} \\
\midrule
LeNet-5 (MNIST) & 0.1ms & 2秒 & 20000× \\
ResNet-20 (CIFAR) & 1ms & 5分钟 & 300000× \\
BERT (小型) & 10ms & 数小时 & $>$100000× \\
\bottomrule
\end{tabular}
\end{table}

\textbf{结论}：FHE目前仅适用于延迟不敏感、高安全需求场景（如医疗诊断、金融审计）。

\subsection{深度技术详解：可信执行环境（TEE）}

可信执行环境提供硬件级数据保护\cite{IntelSGX2024, AMDSEV2024, ARMTrustZone2024}。

\textbf{1. Intel SGX（Software Guard Extensions）}

架构特点：
Enclave：隔离的安全内存区域，CPU直接加密；EPC（Enclave Page Cache）：专用安全内存，当前限制128MB-1GB；远程证明（Remote Attestation）：验证代码完整性和TEE真实性；密封存储（Sealing）：数据持久化保护。

性能特点：
内存加密开销：约5-10\%；Enclave进出开销：约8000-20000 CPU周期；EPC换页开销：严重时性能下降10-100×。

\textbf{2. AMD SEV（Secure Encrypted Virtualization）}

架构特点：
VM级别隔离：整个虚拟机内存加密；SEV-SNP：增加内存完整性保护；支持更大内存：无EPC限制；云友好：适合confidential computing。

\textbf{3. ARM TrustZone/CCA}

架构特点：
双世界架构：Normal World + Secure World；CCA（Confidential Compute Architecture）：Armv9新特性；Realm：类似SGX的隔离区域；适用：移动设备、边缘计算。

\textbf{4. NVIDIA Confidential Computing}

GPU TEE for AI：
H100支持：Confidential Computing模式；显存加密：AES-256全显存加密；Attestation：GPU远程证明；应用：保护推理请求和模型权重。

\textbf{系统架构}：
首先，客户端——验证服务器TEE证明；建立安全通道（TLS + Attestation）；发送加密推理请求。其次，TEE服务器——在Enclave内加载模型权重；解密用户请求；执行推理计算；加密返回结果。第三，信任保证——模型提供方：模型不被服务器运营方窃取；用户：请求内容不被服务器看到；服务器：只需信任Intel/AMD/ARM。

\textbf{性能基准（H100 Confidential模式）}：
吞吐量下降：约15-25\%；延迟增加：约10-20\%；显存开销：加密元数据约5\%；适用场景：高安全敏感推理服务。

\subsection{数据跨境流动技术管控方案}

\textbf{1. 分级分类管控}

\begin{table}[H]
\centering
\small
\begin{tabular}{lccc}
\toprule
\textbf{数据类别} & \textbf{出境要求} & \textbf{技术措施} & \textbf{监管方式} \\
\midrule
核心数据 & 禁止出境 & 物理隔离+加密存储 & 专项审计 \\
重要数据 & 安全评估 & 加密传输+审计日志 & 事前评估 \\
一般数据 & 备案制 & 加密+脱敏可选 & 事后抽查 \\
公开数据 & 自由流动 & 无强制要求 & 无 \\
\bottomrule
\end{tabular}
\end{table}

\textbf{数据出境网关}的设计体现了技术管控与业务效率的平衡。系统首先通过数据分类引擎对出境数据进行自动识别，该引擎结合NLP技术和预设规则实现智能分类。随后，敏感检测模块对数据进行深度扫描，识别身份证号、手机号、地址等个人身份信息（PII）。对于需要脱敏处理的数据，系统采用K-匿名、L-多样性、差分隐私等技术进行自动脱敏。传输控制模块根据分类和脱敏结果决定数据是否允许出境。整个过程由审计留痕模块完整记录，确保所有跨境数据流可追溯、可审计。

\textbf{3. 差分隐私脱敏算法}

对于聚合统计数据出境，使用差分隐私：
\begin{equation}
\tilde{f}(D) = f(D) + \text{Lap}(\Delta f / \varepsilon)
\end{equation}

其中$\Delta f$为全局敏感度，$\varepsilon$为隐私预算。

典型配置：
统计查询：$\varepsilon=1.0$，误差约$\pm 5\%$；机器学习：$\varepsilon=8.0$，模型精度下降约2-5\%；高敏感场景：$\varepsilon=0.1$，误差约$\pm 30\%$。

\subsection{国际格局深度分析}

当前全球AI发展呈现出多极竞争的格局，各主要力量在技术实力、政策取向和战略目标上存在显著差异。

美国在这一领域保持全面领先地位，同时持续强化对竞争对手的遏制。在技术层面，GPT-5系列、Claude 4系列等前沿模型确保了美国的技术优势。在政策层面，芯片出口管制措施持续升级，技术封锁的范围不断扩大。军事化进程明显加速，国防部AI中心已全面运作。在国际治理领域，美国试图主导AI治理规则的制定，塑造有利于自身的国际秩序。

欧盟呈现出监管先行、发展滞后的特征。《AI法案》的全面实施使欧盟在法规层面走在前列，但技术发展相对落后，缺乏世界级的基础模型。这种状况引发了战略焦虑，欧盟担心成为中美技术竞争的"殖民地"，因此在中美之间寻求平衡，试图找到自己的战略空间。

中国正处于快速追赶并谋求自主可控的阶段。DeepSeek-V3等模型已达到世界先进水平，技术追赶的步伐明显加快。尽管芯片禁令造成了一定影响，算力资源受到制约，但中国在部分垂直领域的应用走在前列，并在AI治理方面形成了独特的探索路径。

其他力量同样在全球AI格局中扮演重要角色。英国发挥着"桥梁"作用，积极推动国际对话与合作。日韩在芯片供应链中占据关键位置，是不可忽视的重要参与者。全球南方国家则在争取AI发展权的同时，反对技术霸权，成为塑造国际治理格局的新兴力量。

\subsection{中美博弈的具体策略}

中美AI竞争是长期的、结构性的，但并非零和博弈。在竞争中寻找合作空间，是最优的战略选择。这需要在不同领域采取差异化的策略。

在核心竞争领域必须寸步不让。芯片供应链的自主替代和封锁突破是首要任务。基础模型能力的提升需要持续投入，追赶并力争超越国际前沿。国际标准制定中应积极参与，扩大中国的话语权。人才争夺同样不可放松，既要提升国内环境的吸引力，也要减少高端人才的流失。

在特定领域应积极推动合作。AI安全对齐是双方的共同利益所在，具有合作的天然基础。防止AI在核安全、生物安全等领域的滥用是国际社会的共识，可以成为合作的切入点。学术交流渠道需要保持，这对于跟踪前沿进展和培养人才至关重要。技术标准的互认同样值得探索，避免标准割裂带来的效率损失。

在敏感领域需要坚守底线思维。供应链安全方面必须做好最坏情况的准备，建立足够的战略储备和替代方案。人才安全需要重点关注，防止核心人才被挖角流失。数据安全的防护要到位，防止关键数据被非法获取。舆论安全同样不可忽视，需要有效应对认知战威胁。

\subsection{中欧合作策略}

中欧AI合作具有坚实的现实基础。欧盟在AI治理方面积累了先行经验，其规制框架值得研究借鉴。欧盟对美国技术霸权存在担忧，客观上需要引入制衡力量。中欧在AI伦理理念上存在共同点，对人权保护、隐私权等议题有相近的关切。此外，中欧经济利益存在互补性，合作空间广阔。

合作的重点应聚焦于四个方面。治理对话机制的建立是基础，应推动建立中欧AI治理定期对话机制，增进相互理解和信任。标准互认可以成为实质性合作的突破口，推动AI安全标准的相互认可，降低技术壁垒。学术交流是长期合作的纽带，应扩大AI研究人员的交流往来，促进知识共享。企业合作则是经济利益的交汇点，应支持中欧企业在AI领域开展务实合作。

\subsection{全球南方策略}

全球南方国家的战略意义日益凸显。这些国家是AI发展的潜力市场，人口规模和经济增长带来巨大的应用需求。同时，它们在国际治理中构成重要的力量板块，是塑造全球AI治理格局不可忽视的参与者。

与全球南方国家的合作应采取多元化策略。技术援助是建立长期关系的重要途径，帮助发展中国家建设AI基础设施，弥合数字鸿沟。人才培养可以产生深远影响，通过提供AI培训项目，培育当地的技术生态。标准输出有助于扩大中国技术方案的影响力，推广中国的AI标准和解决方案。在国际治理层面，应与全球南方国家形成话语联合，共同发出发展中国家的声音。

\subsection{日韩合作策略}

日韩在AI领域处于中美之间的微妙位置，这种复杂的三角关系既带来挑战也蕴含机遇。两国既是美国的盟友，又是中国重要的经贸伙伴，在芯片供应链中扮演着关键角色。

对日合作存在现实机遇。日本在机器人、自动驾驶等领域有深厚积累，技术互补性强。老龄化社会对AI应用有强烈需求，市场空间广阔。部分企业对中国市场存在依赖，有合作的经济动力。合作重点可以聚焦于产业应用领域，特别是制造业AI和养老AI，学术研究交流涵盖AI安全和机器人等方向，标准协调则致力于避免技术标准的割裂。需要注意的是，受日美同盟的限制，核心技术合作空间有限。软银深度参与美国Stargate项目，这一动向值得密切关注。

韩国的战略价值主要体现在半导体领域。三星、SK海力士是全球存储芯片的巨头，韩国在半导体设备和材料领域也具备一定能力。合作重点可以包括在合规范围内开展半导体供应链对话，在文娱、游戏、电商等领域推进AI应用合作，以及扩大学术和人才交流。需要提示的风险是，韩国已加入美国芯片联盟，所有合作都需要在合规框架内审慎推进。

\subsection{中东地区合作策略}

中东地区正在积极发展AI产业，拥有充裕的主权财富基金，已成为全球AI投资的新兴力量，具有重要的战略价值。

阿联酋是中东AI发展的领头羊。MGX主权基金参与美国Stargate项目，展现了雄厚的资本实力。技术2071愿景将AI作为国家战略的核心组成部分，发展决心坚定。G42公司是中东领先的AI企业，与中国有合作历史，为进一步合作奠定了基础。合作重点可以包括在第三国设立联合AI投资基金，为阿联酋培训AI人才，以及在智慧城市和能源AI领域开展应用合作。需要注意的是，美国对G42与中国合作施加了压力，相关合作需要谨慎评估风险。

沙特阿拉伯同样蕴含重要合作机遇。沙特2030愿景强调科技转型，AI是其中的重点领域。公共投资基金资金充裕，对外投资活跃。沙特对中国技术相对开放，合作环境相对友好。合作重点可以聚焦于NEOM项目的智慧城市建设、石油化工领域的AI应用，以及中文-阿拉伯语AI模型的联合开发。

以色列具有特殊的战略价值。以色列在AI基础研究方面处于世界领先水平，网络安全AI能力尤为突出，创业生态高度活跃。考虑到美以关系的紧密程度，合作策略应以保持学术层面交流为主，通过第三方渠道了解技术趋势，并积极吸引以色列华人AI人才回流。深度技术合作的空间受到客观限制。

\subsection{东南亚地区合作策略}

东南亚是中国"数字丝绸之路"的重要节点，是AI应用的广阔市场，也是产业转移的潜在目的地，具有重要的战略意义。

合作重点应根据各国特点差异化布局。新加坡作为区域AI枢纽，适合开展高端研发合作和人才交流。印尼、越南、泰国等国家可以重点推进AI应用合作，涵盖电商、物流、金融科技等领域，同时开展数据中心建设和人才培训项目。马来西亚可以借助华人社区的纽带作用，推进中文AI应用的落地。

合作模式应注重系统性和可持续性。"一带一路"框架下的数字基础设施建设是合作的重要载体。设立区域AI培训中心有助于培育当地人才生态。支持当地AI创业生态可以形成长期合作伙伴关系。推广中国AI标准和解决方案则是扩大技术影响力的有效途径。

\subsection{国际话语权建设}

国际话语权的建设需要从框架构建开始，形成系统性的叙事体系。

核心叙事应围绕几个关键主张展开。AI应该造福全人类，而不是成为少数国家的霸权工具。AI发展权是发展中国家的正当权利，不应被人为剥夺。AI治理应该在联合国框架下进行，不能由少数国家垄断规则制定权。技术封锁和脱钩违背科技发展规律，损害全球共同利益。

在此基础上需要提炼几个关键概念。"AI发展权"对标人权话语，强调发展中国家发展AI的正当权利。"技术公平"反对技术霸权和歧视性政策，倡导技术资源的公平分配。"包容性治理"强调多元参与和民主决策，反对少数国家主导。"负责任创新"强调安全与发展并重，追求技术进步与风险防控的平衡。

传播渠道的建设同样重要。国际组织是重要的发声平台，应充分利用联合国、G20等场合阐述立场。学术发表是影响专业话语的有效途径，应积极在顶级会议和期刊上发表研究成果。智库交流有助于影响政策圈层，应与国际智库建立常态化对话机制。媒体传播直接面向公众，应构建多语种对外传播矩阵。

\section{应急准备：最坏情况预案}

\subsection{情景推演}

战略规划必须考虑最坏情况，为极端场景做好准备。以下是三种需要重点防范的情景及其应对预案。

如果美国将对华芯片禁令扩展到所有先进制程芯片，包括从第三国转售的渠道，这将构成芯片全面断供的情景。影响评估显示，短期内一年以内训练新模型的能力将下降50\%以上，中期两到三年内需要依靠存量芯片和国产替代勉强维持，长期影响则取决于国产芯片的突破进度。应对预案应包括保持六到十二个月关键芯片库存的战略储备，最大化现有芯片利用效率的算力优化措施，通过MoE等架构降低算力需求的技术创新，以及国产芯片优先保障关键任务的替代加速方案。

如果美国禁止所有AI技术向中国转移，包括开源模型、学术交流和人才流动，这将导致技术全面脱钩。这种情况的主要影响是丧失跟踪前沿进展的主要渠道，存在形成技术"黑洞"的风险，即不了解对手的技术动向。应对预案应包括建立多元信息渠道，不依赖单一来源，加强自主创新能力以减少外部依赖，扩展与欧洲和全球南方的技术交流，以及支持海外华人学者建立非正式交流网络。

如果发生重大AI安全事件，例如模型被攻破导致敏感信息泄露，需要启动应急响应流程。事件发生后应立即报告国家AI应急响应中心。两小时内完成初步评估，确定影响范围和严重程度。四小时内开始应急处置，隔离受影响系统并启动备份。二十四小时内开展深入调查，查明原因并修复漏洞。七天内完成总结复盘，发布调查报告并制定改进措施。

\section{本章结论}

本章的分析可以提炼为以下几点核心认识。

人才危机是当前最紧迫的瓶颈，我国AI安全人才缺口超过97\%，这一短板如不能尽快补齐，将严重制约整体战略的实施。人才培养需要分类施策，对齐研究者、红队测试员、安全工程师、治理研究者、战略人才这五类人才应根据各自特点采取差异化培养路径。在薪酬待遇上，顶尖人才的薪酬必须具有国际竞争力，建议最高可达300至800万元每年，以有效应对国际人才竞争。

治理架构应采用"一委加一办加多中心"的模式，实现统一协调和专业支撑的有机结合。国际策略需要针对不同对象采取差异化方针：对美"竞合并行"，在竞争中寻求合作空间；对欧"深化合作"，把握治理对话和标准互认的机遇；对全球南方"技术援助加话语联合"，扩大国际支持基础。话语框架的构建应以"AI发展权"、"技术公平"、"包容性治理"为核心，形成系统性的国际话语体系。应急准备方面需要做好芯片断供、技术脱钩、安全事件等最坏情况的预案，确保在极端情况下仍能维持核心能力。

